%% annex/p04/machines-effective-procedures.tex --- proof migrations for
%% v1.1 Part V, ch:turing-machines.  One 40--60 percent compression pass;
%% source-complete items remain complete and source sketches remain sketches.

\begin{annexitem}{annex:p04-machines.1}{thm:p04-pr-functions-computable}{Primitive recursive functions are computable}
\sketchnote{Concrete transition tables and tape-layout invariants are omitted.}
Initial functions have direct machines, and composition sequences retained-input
simulations.  For recursion, store \(\vec x\), \(i\), and the current value;
initialize with \(g\) and run \(h\) exactly \(n\) times.  This finite loop
halts, so induction on generation proves total computability.
\end{annexitem}

\begin{annexitem}{annex:p04-machines.2}{thm:p04-machine-decidability}{Machine form of decidability}
\begin{proof}
A machine computing \(\operatorname{char}_A\) decides by accepting on \(1\) and rejecting on
\(0\); totality gives halting.  Conversely, a decider followed by writing its
Boolean answer computes \(\operatorname{char}_A\).
\end{proof}
\end{annexitem}

\begin{annexitem}{annex:p04-machines.3}{thm:p04-ce-equivalences}{Equivalent forms of computable enumerability}
\sketchnote{Scheduling and history checks are omitted.}
Domain simulation recognizes; dovetailing recognized inputs yields a total
range enumerator, repeating values when nonempty.  Finite accepting histories
are decidable witnesses; witness search recognizes.  Total functions have
nonempty range.
\end{annexitem}

\begin{annexitem}{annex:p04-machines.4}{thm:p04-ce-complement}{Complement criterion for decidability}
\begin{proof}
Swap a decider's answers to recognize the complement.  Conversely dovetail
both recognizers; exactly one accepts, and its identity decides membership.
\end{proof}
\end{annexitem}
